\subsection{Follow-ups: who recommends the brand, and where}
\label{sec:attribution}

Third-party endorsement had the largest effect in the main design, but it added a recommendation, another mention of
the brand and another source at once; the exploratory follow-up in \cref{sec:explore} separates the first two. A second
follow-up separated the wording of a recommendation from its stated source. The author fixed its protocol immediately
before execution on 10~November~2026; the first query ran at 14:24:35~IST, and the public record of the protocol was
committed after the run (\cref{app:prereg}). It reused the 90 trials of the main design and changed only the end of the
target's own page, which received either nothing (a fresh baseline) or one sentence: neutral (``\ldots{} is one option
considered for'' the scenario's stated need) or recommending (``\ldots{} is the option recommended for'' that need),
attributed either to ``the brand's review team'' or to ``the Neralin review team'', a fictional outside reviewer named
nowhere else. Within a trial the four sentences had the same length, need and position and differed only in stance and
source. Of the 5{,}400 planned calls, \AttrNRows{} were made; \AttrNFailedW{} failed with server errors or timeouts and,
as the protocol required, were not retried. Nemotron 3 Super accounted for \AttrNFailedExcludedW{} of them and then
stopped at its spending cap with fewer than the pre-specified \AttrMinTrials{} complete trials, so the analysis covers
\AttrNModelsW{} models and \AttrBlocks{} complete model--trial blocks over 30 scenarios. One answer, from Mistral
Large 3, reached the 4{,}000-token limit and was scored as written.

\begin{table}[tp]
\centering\small
\caption{Attribution follow-up: paired effects in percentage points, pooled over \AttrNModels{} models with equal
weight, with 95\% scenario-cluster bootstrap intervals and Holm correction over the three contrasts within each
outcome. Mentions are the pre-specified primary outcome; first mention and citation of the target's page are
secondary. $^{*}$Holm-adjusted $p<0.05$.}
\label{tab:attribution}
% Generated by make_numbers.py; do not edit.
\begin{tabular}{@{}lrrrr@{}}
\toprule
 & \multicolumn{2}{c}{Target mentioned} & \multicolumn{2}{c}{Named first} \\
\cmidrule(lr){2-3}\cmidrule(l){4-5}
Contrast & Effect [95\% CI] & Holm $p$ & Effect [95\% CI] & Holm $p$ \\
\midrule
Recommending vs neutral wording, brand as source & $+10.1$$^{*}$ {\footnotesize$[6.7, 13.8]$} & <0.001 & $+23.3$$^{*}$ {\footnotesize$[17.9, 29.4]$} & <0.001 \\
Outside reviewer vs brand, recommending wording & $+3.1$$^{*}$ {\footnotesize$[1.4, 4.9]$} & 0.003 & $+4.9$$^{*}$ {\footnotesize$[2.9, 6.8]$} & <0.001 \\
Difference between the two sources' wording effects & $+2.0$ {\footnotesize$[-0.6, 4.6]$} & 0.150 & $+3.2$$^{*}$ {\footnotesize$[1.2, 5.4]$} & 0.008 \\
\bottomrule
\end{tabular}

\end{table}

Weighting models equally, the target was named in \AttrRateB\% of answers with no added sentence, \AttrRateSZero\% and
\AttrRateSOne\% with the brand's neutral and recommending sentences, and \AttrRateQZero\% and \AttrRateQOne\% with the
outside reviewer's. When the brand made the statement, recommending wording raised mentions by \AttrRecU{} points over
neutral wording (95\% CI \AttrRecCI; \cref{tab:attribution}) and the chance of being named first by \AttrTopRecU{}
points. Attributing the same recommendation to the outside reviewer added \AttrExtU{} points to mentions (\AttrExtCI)
and \AttrTopExtU{} to first mentions. Recommending wording raised mentions in all \AttrNModelsW{} models, and the outside
source raised them in \AttrExtPosW{} and lowered them in none; with the brand recommending itself, though,
\AttrSOneHighW{} models already named it in at least 85\% of answers, which leaves little room. Whether the outside
source changed the effect of recommending wording was not resolved for mentions (\AttrIntU{} points, \AttrIntCI). For
first mentions, a secondary outcome, it strengthened it (\AttrTopIntU{} points, \AttrTopIntCI), which does not resolve
the primary contrast. Citation of the target's page, the other secondary outcome, rose with recommending wording by
\AttrCitRecU{} points (\AttrCitRecCI), but the outside source's \AttrCitExtU{} points (\AttrCitExtCI) did not survive
Holm correction. Equal developer weighting gave the same verdicts for mentions. Descriptively, the neutral sentence
alone raised mentions by \AttrSZeroVBU{} points over the fresh baseline: it repeats the target, ties it to the question's
need and adds text, and these cannot be separated. The fresh baseline was close to the original one: over the same
\AttrBasePairs{} trials, \AttrBaseNew\% of its answers named the target, against \AttrBaseOld\% four to six days earlier,
and the two answers agreed on this in \AttrBaseAgree\% of trials. This is a descriptive comparison, not a test of
equivalence.

On these supplied synthetic pages, then, recommending rather than neutral wording raised mentions more than moving the
stated source of the recommendation from the brand's team to a fictional outside team, although the outside source
still raised them. The design tests two wording templates and one invented reviewer, not provenance in general,
independence or credibility; its recommending wording (``the option'') also implies exclusivity; and its sentences sit
on the brand's own page, so it does not test whether the same words carry more weight on a separate site.

\paragraph{A separate site.} A third follow-up tested that. It was specified after the second follow-up's results and
fixed immediately before execution later the same day; the first query ran at 17:31:16~IST, and the public record was
again committed after the run (\cref{app:prereg}). It reused the 90 trials and 12 models with five arms: a fresh
baseline ($B$); the outside reviewer's neutral or recommending sentence at the end of the target's page, word for word
as above ($Q_0$, $Q_1$); and the identical sentence at the end of the buyer's guide, a page on another domain in the
same set, with the target's page left unchanged ($T_0$, $T_1$). On the guide the sentence also puts the target's name
in a second document, which is part of what a separate site is. Three primary contrasts formed one Holm family,
$T_1-Q_1$, $T_0-Q_0$ and $(T_1-T_0)-(Q_1-Q_0)$; no direction was predicted. All \SiteNRows{} calls succeeded and
no answer reached the token limit, so the analysis covers all \SiteNModelsW{} models and \SiteBlocks{} complete blocks.

\begin{table}[tp]
\centering\small
\caption{Separate-site follow-up: the outside reviewer's sentence at the end of the buyer's guide against the same
sentence at the end of the brand's page, in percentage points, pooled over \SiteNModels{} models with equal weight, with
95\% scenario-cluster bootstrap intervals and Holm correction over the three contrasts within each outcome. Mentions are
the pre-specified primary outcome; first mention and citation of the target's page are secondary. $^{*}$Holm-adjusted
$p<0.05$.}
\label{tab:site}
% Generated by make_numbers.py; do not edit.
\begin{tabular}{@{}lrrr@{}}
\toprule
Contrast & Effect & 95\% CI & Holm $p$ \\
\midrule
\multicolumn{4}{@{}l}{\emph{Target mentioned (primary)}} \\
\quad Recommending sentence on the guide vs the brand page & $+0.8$ & $[-0.8, 2.7]$ & 0.441 \\
\quad Neutral sentence on the guide vs the brand page & $+4.6$$^{*}$ & $[1.9, 7.4]$ & 0.013 \\
\quad Difference between the two pages' wording effects & $-3.8$$^{*}$ & $[-6.7, -1.1]$ & 0.031 \\
\multicolumn{4}{@{}l}{\emph{Named first (secondary)}} \\
\quad Recommending sentence on the guide vs the brand page & $+3.1$ & $[-0.5, 6.9]$ & 0.271 \\
\quad Neutral sentence on the guide vs the brand page & $+3.5$$^{*}$ & $[1.5, 5.5]$ & 0.004 \\
\quad Difference between the two pages' wording effects & $-0.5$ & $[-3.7, 3.1]$ & 0.838 \\
\multicolumn{4}{@{}l}{\emph{Target page cited (secondary)}} \\
\quad Recommending sentence on the guide vs the brand page & $-7.7$$^{*}$ & $[-10.2, -5.3]$ & <0.001 \\
\quad Neutral sentence on the guide vs the brand page & $-3.2$ & $[-6.7, 0.1]$ & 0.082 \\
\quad Difference between the two pages' wording effects & $-4.4$ & $[-8.1, -1.0]$ & 0.057 \\
\bottomrule
\end{tabular}

\end{table}

Weighting models equally, the target was named in \SiteRateB\% of answers with no added sentence, \SiteRateQZero\% and
\SiteRateQOne\% with the reviewer's neutral and recommending sentences on its own page, and \SiteRateTZero\% and
\SiteRateTOne\% with the same sentences on the guide. Moved to the guide, the neutral sentence raised mentions by
\SiteNeuU{} points (\SiteNeuCI; \cref{tab:site}), with a positive estimate in \SiteNeuPos{} of the 12 models, and first
mentions by \SiteTopNeuU{} points. For the recommending sentence, close to the ceiling in both places, the difference in
mentions was not resolved (\SiteRec{} points, \SiteRecCI). The third contrast equals $(T_1-Q_1)-(T_0-Q_0)$:
recommending rather than neutral wording added \SiteWordQ{} points on the brand page and \SiteWordT{} on the guide, so
the contrast was \SiteInt{} points (\SiteIntCI). With developers weighted equally it did not survive Holm correction
($p=\SiteDevIntP$), whereas the neutral sentence's gain did, and the ceiling is a plausible reading of it rather than an
identified mechanism. On the guide, the recommending sentence also lowered citations of the target's own page by
\SiteCitRecU{} points (\SiteCitRecCI); descriptively, the guide itself was cited in \SiteGuideCitTOne\% of these
answers, against \SiteGuideCitQOne\% when the sentence was on the brand page, and both pages were cited in
\SiteBothCitTOne\% and \SiteBothCitQOne\% respectively. The arms repeated from the second follow-up came out close
to it: over the \SiteAcrossNW{} models analysed in both, the baseline and the brand-page sentences named the target in
\SiteAcrossB, \SiteAcrossQZero{} and \SiteAcrossQOne\% of answers, against \SiteEarlierB, \SiteEarlierQZero{} and
\SiteEarlierQOne\% earlier.

On these synthetic pages, then, the outside reviewer's neutral sentence named the target more often from the buyer's
guide than from the brand's own page; for its recommending sentence no difference in mentions was detected, while
citations of the brand's page fell and, descriptively, citations of the guide rose. The design tests one synthetic guide and one invented reviewer, not
real independence or reputation, and on the guide the sentence also repeats the target's name in a second document, so
the gain cannot be assigned to the page's independence rather than to that repetition.
